\section{Discussion}
\label{sec:discussion}

\subsection{Why AI Systems Get the Tier Order Wrong}

The training distribution of large language models encodes domain identity as the primary
differentiator among creative professionals. Legal documents in the training corpus are
written by lawyers, annotated as legal, and evaluated against legal genre conventions.
Medical records are written by clinicians and evaluated against clinical documentation
standards. Marketing copy is written by marketers and evaluated against brand and
conversion metrics. When a system learns from this distribution, it learns that the
salient variable distinguishing a lawyer from a doctor from a marketer is their domain
--- and that activating the right domain identity produces domain-appropriate outputs.
This is a locally correct inference about the training data and a systematically wrong
inference about creative quality.

The training data is labeled by domain, but creative quality is defined by receiver
experience. The lawyer's document that a client cannot understand is a domain-correct
creative failure. The marketer's copy that fails to move the reader is a brand-consistent
creative failure. The medical record that the patient cannot act on is a
clinically-accurate creative failure. These failures are invisible in the training
distribution because they are rarely labeled as such: the signal about receiver experience
is sparse, indirect, and unevenly distributed across the corpus. The signal about domain
identity is dense, explicit, and systematic. AI systems learn the available signal, which
is Tier 3, and treat it as primary.

The OLE formula is a prompt-level correction for this training distribution bias. It does
not require retraining; it requires reordering. By establishing the orientation sentence
before any domain content enters the prompt, the practitioner provides a signal that the
training distribution undersupplied: here is who this is for and what they need to
experience. The discipline sentences then operationalize that orientation into the kinds
of checkable constraints that the training distribution does encode --- quality-control
conventions, fact-checking norms, editorial standards --- but directs them at the
receiver's experience rather than the domain's formal expectations.

\subsection{The Orientation as the Hardest Thing to Train}

Tier 1 orientation requires what developmental psychologists call a theory of the other:
the ability to model what another person perceives, understands, and experiences before
making any design decision \cite{premack_woodruff}. In human creative practice, this
capacity is developed through years of training that is explicitly organized around
receiver feedback: the artist shows work and watches how the audience responds; the
designer conducts usability testing; the teacher observes whether students can apply the
material. The feedback loop is direct and constitutive of the training. Artistic education
is, in substantial part, an institutionalized practice for building the theory of the
other that Tier 1 requires.

Language models are trained on the outputs of this process --- on artifacts produced by
people who have already internalized the orientation --- but not on the feedback loops
that produced the orientation itself. The resulting latent capacity is real but requires
activation. The orientation sentence (``Think first about how the other person will
experience this'') is the activation key: it makes the criterion explicit at the point
where the system begins generating, before domain conventions have organized the output
around Tier 3 categories. The single sentence outperforms eleven domain principles
precisely because it activates the capacity that is most latent and most consequential,
rather than the capacity that is most explicit in the training signal.

The mechanism by which the orientation sentence produces its effect remains unestablished;
three candidate accounts are consistent with the data. Under a \emph{distribution-shift
account}, the sentence narrows the generation distribution toward receiver-experience
framings, suppressing domain-convention patterns. Under an \emph{attractor-competition
account}, orientation and domain-convention representations compete in early layers and
the orientation sentence biases this competition. Under a \emph{saliency-weighting
account}, the sentence increases attention to experience-relevant task features,
disproportionately represented in output. These accounts predict different intervention
points and different transfer properties; distinguishing them requires interpretability
methods beyond this study's scope. The computational mechanism of orientation is an open
question.

\subsection{The Symmetric Architecture}

The same Tier 1 orientation that makes a good author makes a good reviewer, because both
roles are ultimately evaluated by the same criterion: the quality of the downstream
person's experience. A reviewer operating from Tier 3 --- surfacing domain-correct
observations without an orientation toward the solver or the author --- produces feedback
that is formally accurate and experientially irrelevant. The author cannot use it to
improve the puzzle; the solver cannot use it to navigate the puzzle. C5's failure in the
reviewing study (lens without philosophy = 21.7, below the 23.7 bare baseline) is the
reviewing-side analog of the authoring-side L-T1 failure (experience-first removed =
12.3, categorical failure). Both failures occur for the same structural reason: Tier 3
without Tier 1 produces outputs organized around domain format rather than receiver
experience.

This symmetry has implications beyond the two roles studied here. Any role in a creative
production chain --- editing, curating, directing, commissioning --- that involves
evaluating or shaping creative work for a downstream receiver should exhibit the same
tier structure. The Tier 1 orientation, expressed appropriately for each role, is the
universal entry point for quality in the chain.

This study has three limitations that bound the generalization. First, the empirical base
is a single creative domain: puzzle-hunt design, which has unusual properties (fully
verifiable solutions, established evaluation rubrics, a specific solver-experience
criterion) that may not transfer directly to domains where quality criteria are more
contested. Second, the authoring and reviewing agents are AI systems simulating human
expert roles rather than human professionals responding to prompts; the transfer to
human-in-the-loop creative systems requires separate validation. Third, this study
examines only authoring and reviewing roles; other roles in the creative production chain
--- editing, directing, commissioning --- may exhibit different tier weightings or
different orderings within Tier 2.
